% =============================================================
% Subsection: Welfare and heterogeneity (inside Section 4)
% =============================================================

\subsection{Welfare and incidence}\label{sec:welfare}

The welfare calculation is best read as a consumption-leisure diagnostic. It asks a narrow question: if employment is fixed, wages do not bargain, and the reform mechanically gives incumbent workers more non-work time, how large is the leisure gain relative to the consumption loss? It is not a social-welfare theorem. Health, household production, bargaining, prices, employment transitions, and distributional weights are outside the calculation.

Within that narrow object, the schedule has a clear shape. Around 41 to 42 hours, the model buys leisure at a relatively low output cost. Below 40 hours, the price of additional leisure rises sharply because the reform is no longer mostly trimming long-hours fatigue. That is why the representative-agent welfare measure turns negative before the endpoint. The one-sided fatigue and two-sided penalty specifications agree on the moderate-reduction region and disagree mainly once the cap moves below the efficiency peak.

\begin{figure}[htbp]
\centering
\includegraphics[width=0.92\textwidth]{fig_welfare_schedule.pdf}
\caption{Welfare schedule by hours cap under both calibrations.}\label{fig:welfare_schedule}
\vspace{2pt}
\begin{minipage}{0.92\textwidth}
\justifying
{\footnotesize \textit{Notes:} Each line plots aggregate welfare $\Delta$CV from Equation~\eqref{eq:cv} at successively lower hours caps, relative to the $44$h baseline, under the one-sided fatigue and two-sided penalty $e(h)$ specifications. $\Delta$CV peaks near $42$h at about $+0.62\%$ and turns negative below $40$h. At the $36$h endpoint, welfare falls by $1.76\%$ under the one-sided specification and by $3.63\%$ under the two-sided one.}
\end{minipage}
\end{figure}

A TFP gain of approximately $1.4\%$ at the 36h cap is enough to make the endpoint welfare-neutral under the one-sided fatigue specification. The two-sided threshold is $2.9\%$. Both thresholds lie below the output-neutral $A_{\text{req}}$ because leisure partially compensates the consumption loss. Appendix Table~\ref{tab:welfare_preference_robustness} shows that this qualitative ranking survives a separable log and CRRA(2) diagnostic: the 40-hour alternative remains positive, while the 36-hour endpoint remains negative.

The aggregate welfare number masks incidence. In the conservative type-level decomposition reported in Appendix Table~\ref{tab:welfare_incidence_by_type}, formal stayers gain $+1.41\%$ in first-order compensating variation because they work roughly six fewer hours. Workers pushed into informality lose $-28.57\%$. This group is small, $1.9\%$ of workers, but its loss is large. Informal workers not affected by the statutory cap have zero direct incidence in the model. The representative-agent value in Equation~\eqref{eq:cv} differs from the share-weighted type sum because the GHH composite is nonlinear. For this reason, the 41 to 42h welfare maximum should be read as a model-internal representative-agent result, with distributional incidence reported separately.

The welfare schedule does not imply opposition to reductions in working hours. Under the model's calibration, the welfare-improving margin lies near 41 to 42 hours rather than at 36. The schedule separates output loss, required productivity, and welfare. These objects move differently with the cap, and the welfare-improving region is narrower than the output-neutral one.

Adjustment costs and informality differ across firm sizes and sectors, producing heterogeneous incidence of the reform. Small firms ($\leq 49$ employees) require a larger TFP gain than large firms, reflecting their higher formalization wedges and baseline informality. Applying the national framework separately to three broad sectors (agriculture, industry, services) with sector-specific PNAD parameters delivers a per-sector $A_{\text{req}}$ of $1.3$ to $1.8\%$ for the 40-hour alternative and $6.91$ to $8.05\%$ for the 36-hour endpoint under the two-sided penalty specification. Informality rises in every sector at the 36-hour endpoint (agriculture $+2.88$ pp, industry $+2.73$ pp, services $+2.10$ pp). These are sector-by-sector partial-equilibrium counterfactuals, not an input-output incidence analysis. Industry's $A_{\text{req}}$ slightly exceeds agriculture's because the moderate $\eta_I=0.40$ combined with industry's larger formal base dampens the mechanical gap between the two sectors. Services account for most of the aggregate output loss through their output share, not through a larger per-worker effect. The online appendix reports the full decomposition and the parameter sensitivities.
